\subsection{Particle Propagation}
Following the generation of the primary neutrino interaction via the \gls{marley} generator for low-energy solar neutrino events, the resulting final-state particles are propagated through the liquid argon medium using a \gls{geant4} simulation~\cite{GEANT4_2003}. This so called \emph{g4} stage, models the particle-matter interactions, including energy loss, scattering, bremsstrahlung, photon production, hadronic interactions and more. This happens according to a physics list tailored for these low-energy processes (\href{https://github.com/DUNE/dunesim/blob/develop/dunesim/Simulation/larg4services_dune.fcl}{larg4services}). The result of this first step is a detailed list of daughter particles and their trajectories, along with the energy deposited in the liquid argon at each step of the particle's path.

As these traverse the liquid argon, charged particles will lose energy primarily through ionisation, producing free ionisation electrons and excitation of argon atoms, which subsequently decay and emit prompt scintillation photons in the \gls{vuv} range (\SI{128}{\nano\meter}). These processes are central to the signal generation in a \gls{lartpc}. For computational efficiency, the simulation does not explicitly track every ionisation electron or \gls{vuv} photon. Instead, energy loss along particle trajectories is recorded as a set of energy depositions, or \texttt{edeps}, which include position, time, energy, and associated particle identity. Similarly, the number of scintillation photons produced in each energy deposition step is recorded as a secondary output, commonly referred to as \texttt{simphotons}. These \texttt{edeps} and \texttt{simphotons} are then effectively propagated through the \gls{lar} volume with custom analytical algorithms.

The computation of light and charge production for the \gls{dune} \gls{fd} utilizes the \gls{larql} (Liquid Argon Charge and Light) model for correlated signal generation. \gls{larql} is a phenomenological approach implemented to provide an improved description of the anti-correlation between free ionization charge ($Q$) and scintillation light ($L$). It is particularly effective for \gls{dune} simulations as it covers electric field ranges from \qtyrange{0}{0.75}{\kilo\volt\per\centi\meter} and energy densities ($dE/dx$) from \qtyrange{2}{40}{\mega\electronvolt\per\centi\meter}. The model is built upon the unitary charge-light equation of the traditional Birks' model:

\begin{align}
    Q\left(\dfrac{dE}{dx}, \mathcal{E}\right) + L\left(\dfrac{dE}{dx}, \mathcal{E}\right) = N_i + N_{ex}
\end{align}

where $N_i$ is the number of ionization electrons initially generated ($N_i = 1/W_{ion}$) and $N_{ex}$ is the number of primary excitons.

Considered the evidence of light reduction at $\mathcal{E} \to 0$ for lower $\frac{dE}{dx}$ values~\cite{Doke_2002}, one can establish the fraction of missing photons ($1-\eta_0$) and from this the fraction of escape electrons ($\chi_0$):

\begin{equation}
    \eta_0 = \frac{1 - \chi_0 + N_{ex}/N_i}{1 + N_{ex}/N_i} \quad \Rightarrow \quad \chi_0 = (1 + N_{ex}/N_i) \cdot (1 - \eta_0)
\end{equation}

With this, the \gls{larql} modification to the traditional Birks' model adds a term for these escaping electrons ($\chi_0$):

\begin{equation}
    Q_\text{LArQL} = \frac{A_B}{W_\text{ion}\cdot\left(1 + \frac{k_B}{\mathcal{E}\rho_\text{LAr}} \cdot \frac{dE}{dx}\right)} + \chi_0\left(\dfrac{dE}{dx}\right) \cdot f_{corr}\left(\dfrac{dE}{dx}, \mathcal{E}\right) \cdot Q_\infty
\end{equation}

where $A_B$ and $k_B$ are original Birks parameters, $\rho_{LAr}$ is the \gls{lar} density and $f_{corr}$ is an exponentially decreasing function that reduces the contribution of escape electrons as the field increases which is defined as:

\begin{equation}
    \large
    f_{corr} = e^{-\mathcal{E}\cdot\left(\alpha \ln(\frac{dE}{dx}) + \beta\right)^{-1}}\,.
\end{equation}

Here, the parameters $\alpha$ and $\beta$ are optimized through global fits of experimental data from ICARUS, MicroBooNE, and ARIS.

After the generation of the ionization charge and scintillation light in the \emph{g4} stage, the simulation proceeds to propagate the photons wit a so called \emph{FastSim} approach. The method relies on a \emph{Photon Library} which is pre-computed using a standalone \gls{geant4} simulation that deposits a fixed number of photons in each voxel of the detector volume and records the response (number of photoelectrons) for every photon detector. During the standard simulation, the software looks up the pre-calculated response for the voxel where light was produced, vastly reducing CPU time while still providing an accurate representation of the light collection efficiency across the detector volume~\cite{semiana}.
